\chapter{The Five Diagnostic Questions as a Scientific Instrument}
\index{diagnostic questions}\index{diagnostic audit}

\section{From Theorems to Audits}

The five diagnostic questions of Chapter~4 operate at three levels. As \emph{verification criteria}, they determine whether reconstruction is guaranteed. As a \emph{scientific workflow}, they determine why reconstruction works or fails. As a \emph{research program generator}, they determine what experiment, theorem, or measurement would resolve remaining uncertainty.

The purpose of the audit is not to certify reconstruction. It is to identify the minimum evidence required to justify a reconstruction claim.

\section{The Five-Question Audit Procedure}

\textbf{Question 1: What is the interaction decomposition, and is it exact?}
Specify $\mathcal{X}$, $x_0$, the interaction order, and the terms $I_r^{Q,x_0}$. A positive answer establishes that the query's dependence on the history can be analyzed at all.

\textbf{Question 2: Is the stability sequence summable?}
Specify $\{\sigma_r\}$ and the tail behavior. If not summable: \emph{stability failure} (structural). Report as a property of the trajectory space; consider alternative decompositions.

\textbf{Question 3: Is the compression sufficient?}
Specify $\pi$, $d_\pi$, and the sufficiency error $\varepsilon_1$. If large: \emph{sufficiency failure} (architectural). Redesign the compression.

\textbf{Question 4: Are depths aligned?}
Verify $d_\pi = d_Q$. If not: \emph{misalignment failure}. Adjust depth parameters.

\textbf{Question 5: Does the equivalence relation avoid degeneracy?}
Check whether $\sim_{d_\pi}$ collapses trajectories the query family distinguishes. If so: \emph{degeneracy failure}. Enrich the equivalence relation or measure structure.

\section{The Diagnostic Audit Report}

A complete audit produces a nine-component report:
\begin{enumerate}
  \item Trajectory Space Definition
  \item Query Family Definition
  \item Interaction Decomposition
  \item Stability Analysis
  \item Compression Architecture
  \item Sufficiency Analysis
  \item Alignment Analysis
  \item Failure Classification
  \item Open Empirical Obligations
\end{enumerate}

Component~9 is the most important: it converts vague uncertainty into concrete scientific questions.

\section{Research Programs from Audit Gaps}

Open obligations from the three verified domains:
\begin{itemize}
  \item \textbf{Learning systems:} Measure derivative tensor decay at frontier-model scale.
  \item \textbf{Semantic systems:} Measure attenuation coefficients for natural-language inference grammars.
  \item \textbf{Physical fields:} Determine whether accessible volume contracts exponentially, polynomially, or not at all under RSVP dynamics.
\end{itemize}

The framework transforms philosophical uncertainty about whether reconstruction is possible into empirical questions about specific measurable quantities.

\section{Portability}

The five-question audit is portable to any domain where counterfactual reconstruction from compressed representations is claimed or desired. It is modular, falsifiable, research-program-generating, and sensitive to epistemic levels (proved versus empirically supported versus hypothetical).

The purpose of Part~Three has not been to prove that reconstruction succeeds. It has been to determine what evidence would be required to believe that it succeeds.
